\section*{अध्याय \ref{ch:g7:short-circuits-safety} --- लघुपथन और विद्युत सुरक्षा}

\begin{solution}{exo:g7:short-circuits-safety:1}
एक ही बिंदुओं को जोड़ती दो सड़कों में से धारा उसी पर भीड़ बना लेती है
जो उसका कम विरोध करती है। काम करता लैंप क़तार का विरोध करता है; नंगा
मोटा तार बमुश्किल करता है --- इसलिए क़तार लैंप छोड़कर तार पर चली जाती
है।
\end{solution}

\begin{solution}{exo:g7:short-circuits-safety:2}
उपकरण का \omterm{def:g7:short-circuits-safety:device}{लघुपथन}: किसी एक उपकरण पर बाईपास --- वह काम करना बंद कर देता
है और फंदा चलता रहता है। \omterm{def:g3:first-electric-circuit:battery}{बैटरी} का \omterm{def:g7:short-circuits-safety:device}{लघुपथन}: \omterm{def:g3:first-electric-circuit:battery}{बैटरी} के अपने सिरों को
जोड़ती वह नंगी सड़क जिस पर कोई उपकरण ही न हो --- यह भारी मामला है:
क़तार का विरोध कुछ भी नहीं करता, और भगदड़ तार तथा \omterm{def:g3:first-electric-circuit:battery}{बैटरी} को बरबादी और
आग की ओर तपा देती है।
\end{solution}

\begin{solution}{exo:g7:short-circuits-safety:3}
बाईपास \omterm{def:g3:first-electric-circuit:bulb}{बल्ब} 2 की पूरी क़तार चुरा लेता है --- त्यागा हुआ \omterm{def:g3:first-electric-circuit:bulb}{बल्ब} बुझ जाता
है। एक विरोधी के बाईपास होते ही पूरा फंदा \omterm{def:g3:first-electric-circuit:battery}{बैटरी} का कम विरोध करता है,
क़तार मज़बूत हो जाती है, और अब उस मज़बूत धारा को अकेले ढोता \omterm{def:g3:first-electric-circuit:bulb}{बल्ब} 1
ज़्यादा चमकीला हो जाता है।
\end{solution}

\begin{solution}{exo:g7:short-circuits-safety:4}
चलती हुई धारा अपनी सड़क को हमेशा गरम करती है, और भगदड़ उसे बेतहाशा
गरम कर देती है। लैंप इसी गरमाहट को जान-बूझकर काम में लेता है: उसका
तंतु सफ़ेद-गरम चलने और चमकने के लिए ही बना है।
\end{solution}

\begin{solution}{exo:g7:short-circuits-safety:5}
जान-बूझकर पतला रखा गया वह क़ुरबान होने वाला तार जो धारा के अपनी तय
सीमा पार करते ही पिघल जाता है --- और फंदा काट देता है; वह ख़ुद को
क़ुरबान करता है। \omterm{def:g4:series-circuits:series}{श्रेणी में} इसलिए, कि भगदड़ को उसी नाकाबंदी से गुज़रना
पड़े: \omterm{def:g5:series-parallel-circuits:parallel}{समांतर में} लगा पहरेदार तो उसे बस गुज़रता देखता रहता।
\end{solution}

\begin{solution}{exo:g7:short-circuits-safety:6}
विच्छेदक पिघलने के बजाय खुल जाता है, और एक झटके से फिर चढ़ा दिया जाता
है --- कोई पुरज़ा नहीं बदलना पड़ता। पर किसी भी पहरेदार को दोबारा चढ़ाने
या बदलने से पहले ख़राबी ढूँढ़कर ठीक करनी चाहिए: वह किसी वजह से ही गिरा
था।
\end{solution}

\begin{solution}{exo:g7:short-circuits-safety:7}
अतिभार: तीन बेगुनाह उपकरणों ने एक ही \omterm{def:g5:series-parallel-circuits:parallel}{शाखा} से अपनी तीनों पूरी क़तारें
एक साथ ढोने को कहा, और जुड़ी हुई दौड़ तारबंदी को तपा देती है। विच्छेदक
उस जोड़ पर गिरता है --- और ठीक ही।
\end{solution}

\begin{solution}{exo:g7:short-circuits-safety:8}
कील कभी \omterm{def:g2:ice-water-steam:meltfreeze}{पिघलती} ही नहीं: वह पहरेदार को चुप कर देती है जबकि बीमारी
भड़कती रहती है। अगली भगदड़ को कोई नाकाबंदी नहीं मिलती, और घर की दीवारों
की तारबंदी ही \omterm{def:g7:short-circuits-safety:fuse}{फ़्यूज़} बन जाती है --- वहाँ झुलसती हुई जहाँ कोई देख नहीं
सकता।
\end{solution}

\begin{solution}{exo:g7:short-circuits-safety:9}
उपकरण से पहले हुआ \omterm{def:g7:short-circuits-safety:device}{लघुपथन} --- यानी उस तार की नज़र से आपूर्ति पर ही
लगा शॉर्ट। लैंप मर जाता है (त्यागा हुआ), भगदड़ छुए हुए चालकों में से
दौड़कर केबल तपा देती है --- और \omterm{def:g5:series-parallel-circuits:parallel}{शाखा} का पहरेदार, \omterm{def:g7:short-circuits-safety:fuse}{फ़्यूज़} हो या
विच्छेदक, \omterm{def:g5:series-parallel-circuits:parallel}{शाखा} काट देता है। वह कटाव ही तंत्र का काम करना है।
\end{solution}

\begin{solution}{exo:g7:short-circuits-safety:10}
दोनों बल्बों पर एक साथ रखा तार --- सादे तार और \omterm{def:g3:first-electric-circuit:battery}{बैटरी} की लीड के रास्ते
--- \omterm{def:g3:first-electric-circuit:battery}{बैटरी} के अपने दोनों सिरों को जोड़ देता है, और बाईपास वाली सड़क पर
कोई उपकरण नहीं होता: यानी \omterm{def:g3:first-electric-circuit:battery}{बैटरी} का \omterm{def:g7:short-circuits-safety:device}{लघुपथन}। दोनों \omterm{def:g3:first-electric-circuit:bulb}{बल्ब} मर जाते हैं और
भगदड़ तार तथा \omterm{def:g3:first-electric-circuit:battery}{बैटरी} को तपा देती है; खिलौना-पैमाने पर भी यह सेल बरबाद
करता है और उँगलियाँ गरम कर देता है, इसलिए इसका प्रदर्शन शब्दों में
होता है, कर के नहीं।
\end{solution}

\begin{solution}{exo:g7:short-circuits-safety:11}
मरते हुए \omterm{def:g3:first-electric-circuit:bulb}{बल्ब} के जुड़ गए सहारे के तार भीतर ही बना-बनाया बाईपास बन गए:
फंदा उस मरे हुए \omterm{def:g3:first-electric-circuit:bulb}{बल्ब} के शॉर्ट \emph{में से होकर} बंद बना रहता है,
इसलिए लड़ी उसका नुक़सान झेल जाती है --- और एक विरोधी कम होने से ज़रा
ज़्यादा चमकीली भी। पर ऐसी हर मौत बचे हुए बल्बों में से गुज़रती क़तार
मज़बूत कर देती है और अगली नाकामी जल्दी ले आती है: लड़ी शालीनता से अपनी
ही झड़ी तैयार कर रही है।
\end{solution}

\begin{solution}{exo:g7:short-circuits-safety:12}
साझा सड़क पर एक मुख्य पहरेदार और मुख्य \omterm{ex:g3:first-electric-circuit:switch}{स्विच}; और फिर हर \omterm{def:g5:series-parallel-circuits:parallel}{शाखा} ---
बत्तियाँ, औज़ार, चार्जर --- अपने \omterm{ex:g3:first-electric-circuit:switch}{स्विच} और अपनी \omterm{def:g5:series-parallel-circuits:parallel}{शाखा} के हिसाब से चुने
गए अपने पहरेदार के साथ (औज़ारों की \omterm{def:g5:series-parallel-circuits:parallel}{शाखा} सबसे भारी क़तारें ढोती है और
सबसे मज़बूत विच्छेदक की हक़दार है; चार्जर की पतली \omterm{def:g5:series-parallel-circuits:parallel}{शाखा} को सबसे नाज़ुक
\omterm{def:g7:short-circuits-safety:fuse}{फ़्यूज़} चाहिए, क्योंकि वह उन धाराओं पर झुलस जाती जिन्हें मुख्य पहरेदार
अनदेखा कर देता)। दराज़ के ऊपर की सूचना, जैसे: ``उड़ने वाला \omterm{def:g7:short-circuits-safety:fuse}{फ़्यूज़} सच
बोल रहा है --- ख़राबी ढूँढ़ो। जैसा था वैसा ही लगाओ; पन्नी और कीलें
किसी की रक्षा नहीं करतीं, बस आग का रास्ता खोलती हैं।''
\end{solution}

\begin{solution}{pb:g7:short-circuits-safety:1}
\textbf{1.} अतिभार: गरमी की भूखी तीन चीज़ों ने अपनी ईमानदार क़तारें एक
ही \omterm{def:g5:series-parallel-circuits:parallel}{शाखा} पर जोड़ दीं --- किसी नंगे तार की ज़रूरत ही नहीं; अकेला जोड़ ही
तारबंदी को सुरक्षा से आगे तपा देता है।
\textbf{2.} कुतरी हुई जगह पर \omterm{def:g7:short-circuits-safety:device}{लघुपथन}: उसी पल से क़तार आगे वाले लैंप को
छोड़कर छूते हुए चालकों में से भगदड़ बनकर दौड़ी और तार तपा दिया।
\textbf{3.} उसने सुरक्षा ख़त्म ही कर दी: पन्नी किसी भी भगदड़ को बिना
पिघले ढो ले जाती है --- \omterm{def:g5:series-parallel-circuits:parallel}{शाखा} पहरेदार-रहित खड़ी रह गई।
\textbf{4.} संभव कहानी: पहले अतिभार आया, और असली \omterm{def:g7:short-circuits-safety:fuse}{फ़्यूज़} ने ---
ईमानदारी से --- शायद एक से ज़्यादा बार उड़कर जवाब दिया; किसी ने इस
``झंझट'' का इलाज पन्नी से कर दिया (सबूत तीन); और जब आख़िर कुतरे तार का
\omterm{def:g7:short-circuits-safety:device}{लघुपथन} हुआ, तब कोई पहरेदार बचा ही नहीं था, इसलिए भगदड़ ने उड़े हुए
\omterm{def:g7:short-circuits-safety:fuse}{फ़्यूज़} और अँधेरी रसोई के बजाय दीवार की तारबंदी को आग तक तपा दिया।
\textbf{5.} धारा अपनी सड़क को हमेशा गरम करती है; दौड़ जितनी तेज़,
गरमाहट उतनी उग्र। लैंप साबित करता है कि गरमाहट साधी जा सकती है ---
उसका तंतु सुरक्षित \omterm{def:g3:first-electric-circuit:bulb}{बल्ब} में सफ़ेद-गरम चलने के लिए ही बना है। यहाँ कमी
थी: गरमी के रहने की कोई बनाई हुई जगह, और उसे रोकने वाला कोई पहरेदार।
\textbf{6.} जान-बूझकर पतला, ताकि परिपथ के सारे तारों में से वही सबसे
पहले पिघले: पहरेदार जानलेवा गरमाहट अपने ही तीन \omterm{def:g2:measuring-length:units}{सेंटीमीटर} में ले लेता
है --- ख़ुद मरकर, ताकि दीवारों के \omterm{def:g2:measuring-length:units}{मीटर} बचे रहें।
\textbf{7.} पहरेदार को सही पल पर \emph{नाकाम} होना ही चाहिए: उसका गुण
एक बनाई हुई कमज़ोरी है। पन्नी का \omterm{def:g4:conductors-insulators:conductor}{चालक} के रूप में उत्कृष्ट होना ही
नाकाबंदी का भ्रष्टाचार था --- रास्ता बेहतरीन, सुरक्षा शून्य।
\textbf{8.} सुरक्षा शिष्टाचार की तरह नहीं जुड़ती: एक ही सड़क पर तीन
वैध क़तारें जुड़कर एक अवैध दौड़ बन जाती हैं --- \omterm{def:g5:series-parallel-circuits:parallel}{शाखा} नीयत नहीं,
अंकगणित मानती है।
\textbf{9.} पन्नी। अतिभार और कुतरे तार वे दुर्घटनाएँ हैं जिन्हें
झेलने के लिए तंत्र बना ही है --- पहरेदार ठीक उन्हीं के लिए हैं; पहरेदार
को चुप कराने ने झेली जा सकने वाली ख़राबियों को आपदा में बदल दिया।
\textbf{10.} ढंग का \omterm{def:g7:short-circuits-safety:fuse}{फ़्यूज़} होता तो: तार का \omterm{def:g7:short-circuits-safety:device}{लघुपथन} होता, दौड़ उछलती,
\omterm{def:g7:short-circuits-safety:fuse}{फ़्यूज़} पल भर में पिघल जाता --- रसोई अँधेरी हो जाती, परिवार लौटकर एक
उड़ा हुआ \omterm{def:g7:short-circuits-safety:fuse}{फ़्यूज़} और एक कुतरा हुआ तार पाता, बड़बड़ाता, दोनों ठीक करता,
और कभी जान ही नहीं पाता कि वह रात कितनी बुरी हो सकती थी।
\textbf{11.} गरमी बनाने वाली हर चीज़ को दीवार के अपने सॉकेट पर, अपनी
ठीक से संरक्षित \omterm{def:g5:series-parallel-circuits:parallel}{शाखा} के साथ; और मल्टी-सॉकेट फिर से सिर्फ़ लैंपों और
चार्जरों के लिए, इस पक्के नियम के साथ: एक सॉकेट पर एक भूखा उपकरण, और
कोई पट्टी कभी किसी हीटर को नहीं खिलाएगी।
\textbf{12.} जैसे: ``खलनायक का तरीक़ा है बिना मेहनत वाली सड़क: क़तार
को बाईपास दो और वह ले लेगी, और दौड़ते हुए तपाती जाएगी। पहरेदार की
ड्यूटी है \omterm{def:g4:series-circuits:series}{श्रेणी में} खड़े रहना और सबसे पहले, ईमानदारी से तथा सस्ते में
नाकाम हो जाना। और कोई भी घर पहरेदार अपने मन से नहीं गढ़ सकता: \omterm{def:g7:short-circuits-safety:fuse}{फ़्यूज़}
की जगह \omterm{def:g7:short-circuits-safety:fuse}{फ़्यूज़} ही आता है --- जैसा का तैसा --- वरना किसी दिन पिघलने का
काम दीवारें ख़ुद कर लेंगी।''
\end{solution}
